\documentclass[journal=langd5,manuscript=article]{achemso}

\usepackage{amsmath}
\usepackage{graphicx}
\usepackage{booktabs}
\usepackage{xcolor}
\usepackage{siunitx}
\usepackage{hyperref}

% -----------------------------------------------------------------------
% Authors — fill in before submission
% -----------------------------------------------------------------------
\author{[Author Names]}
\affiliation{[Institution, Department, Address]}
\email{chalakon2550@gmail.com}

\title{Spontaneous Adsorption of \texorpdfstring{$\beta$}{beta}-Lactoglobulin
       at the Air--Water Interface: Pre-Adsorption Conformational Activation
       Revealed by Unbiased Atomistic Molecular Dynamics}

\abbreviations{BLG, $\beta$-lactoglobulin; MD, molecular dynamics;
               SASA, solvent-accessible surface area; RMSD, root-mean-square
               deviation; RMSF, root-mean-square fluctuation;
               PBC, periodic boundary conditions; PME, particle mesh Ewald}

\keywords{beta-lactoglobulin, air-water interface, molecular dynamics simulation,
          protein adsorption, conformational activation, foam stability,
          hydrophobic exposure, CHARMM36m}

% -----------------------------------------------------------------------
\begin{document}
% -----------------------------------------------------------------------

\begin{abstract}
\noindent
$\beta$-Lactoglobulin (BLG), the primary whey protein in bovine milk, stabilizes
foam through adsorption at the air--water interface; however, the atomic-level
mechanism governing this process has remained elusive. Here we report the first
unbiased atomistic molecular dynamics (MD) simulation of BLG at the air--water
interface, comprising a 1{,}000~ns trajectory from bulk solution and three
500~ns near-interface replicas in a 12$\times$12$\times$35~nm slab geometry
using the CHARMM36m force field and TIP3P water. We discover a \emph{pre-adsorption
conformational activation} mechanism: BLG undergoes transient partial unfolding
in bulk water — manifested as dramatic spikes in hydrophobic solvent-accessible
surface area (SASA, up to \SI{52.5}{\nano\metre\squared}, 2.7-fold above
baseline) — that precedes and enables interfacial approach. During these
activation events the beta-barrel is transiently disrupted while the
alpha-helix (residues 130--142) remains structurally intact, suggesting the
helix acts as a scaffold maintaining calyx orientation during insertion.
The 1{,}000~ns simulation from bulk establishes for the first time by atomistic
simulation that spontaneous BLG adsorption requires $>$1~\textmu{}s, consistent
with the experimentally observed kinetic energy barrier. Per-residue flexibility
analysis identifies Loop~BC (residues 30--35, RMSF $\approx$\,0.53~nm) as the
predicted first-contact region. These findings provide a molecular-level
explanation for the kinetic barrier to BLG adsorption and open new avenues for
rational engineering of milk foam stability.
\end{abstract}

% -----------------------------------------------------------------------
\section{Introduction}
% -----------------------------------------------------------------------

The stability of milk foam is central to the texture and sensory quality of
aerated dairy products ranging from cappuccino to whipped cream. At the molecular
level, foam stability is governed by the proteinaceous film that forms at the
air--water interface, where adsorbed proteins create a viscoelastic layer that
resists bubble coalescence and drainage.\cite{Dickinson1999,BertonCarabin2018}
$\beta$-Lactoglobulin (BLG), the most abundant whey protein in bovine milk
($\sim$3~g\,L$^{-1}$), is the dominant contributor to this interfacial
film.\cite{Cornec1999,Bolisetty2022}

Decades of experimental investigation have established the phenomenology of BLG
adsorption at the air--water interface. Tensiometry, neutron reflectometry, and
ellipsometry show that BLG adsorbs over timescales of seconds to minutes,
exhibits a kinetic energy barrier that distinguishes it from more flexible
proteins such as $\alpha$-lactalbumin, and partially unfolds upon adsorption
to form a rigid interfacial layer.\cite{Cornec1999,Ulaganathan2019,Foam4_2020}
The structure of the adsorbed film has been resolved at \r{A}ngstr\"{o}m
resolution by neutron reflectometry, revealing a protein monolayer of
$\sim$3--4~nm thickness whose rheological properties correlate with foam
stability.\cite{Ulaganathan2019,Foam4_2020} Yet despite this rich experimental
picture, the atomic-level mechanism of adsorption --- what happens to BLG in
bulk solution before and during interfacial insertion --- has remained unknown.

Atomistic molecular dynamics simulation offers a complementary view that
experiment cannot yet provide. Prior MD studies of BLG have focused
exclusively on oil--water interfaces, where the protein was pre-positioned
near the interface and simulated for tens to hundreds of nanoseconds:
Jordens \textit{et al.}\cite{Jordens2015} characterised adsorption at a
decane--water interface and found that hydrophobic residue exposure drives
insertion while secondary structure is largely preserved; Zare \textit{et al.}\cite{Zare2016}
extended this to three oils of varying hydrophobicity. For the air--water
interface --- a fundamentally different system with surface tension
$\sim$\SI{72}{\milli\newton\per\metre} versus $\sim$\SI{20}{\milli\newton\per\metre}
for decane --- no atomistic MD simulation of BLG has been reported.
The closest related study, by Chaudhri \textit{et al.},\cite{Chaudhri2024}
examined monoclonal antibody adsorption at the water/vapour interface using
CHARMM36m over 200~ns, but employed thermally pre-stressed protein conformations
rather than native-state spontaneous dynamics.

Here we present the first unbiased atomistic MD simulation of BLG at the
air--water interface. We ran a 1{,}000~ns trajectory with BLG starting from
the centre of a water slab (SET~1A) and three 500~ns replicas with BLG starting
2.183~nm from the upper interface (Plan~B). We discover a pre-adsorption
conformational activation mechanism, establish the $>$1~\textmu{}s spontaneous
adsorption timescale, and identify the structural elements responsible for the
kinetic energy barrier observed experimentally.

% -----------------------------------------------------------------------
\section{Methods}
% -----------------------------------------------------------------------

\subsection{Protein Preparation}

The crystal structure of bovine BLG (PDB: 1BEB, chain~A, X-ray at 1.8~\r{A}
resolution)\cite{Brownlow1997} was used as the starting configuration. The
monomeric form was selected as the physiologically relevant species at neutral
pH and low protein concentration. Protonation states were assigned at pH~7.0
using the standard CHARMM36m protocol. Sodium and chloride ions were added to
neutralise the system charge.

\subsection{Slab System Construction}

The simulation box measured 12~$\times$~12~$\times$~35~nm, configured as a
slab geometry with a $\sim$7~nm water slab (TIP3P\cite{Jorgensen1983}) centred
at $Z \approx 17.5$~nm, flanked by $\sim$14~nm of vacuum ($\sim$7~nm on each
side). This geometry creates two equivalent air--water interfaces without
artificial constraints.

For SET~1A, the protein was placed at the centre of the water slab
($Z \approx 17.5$~nm). For the Plan~B replicas (1B-R1, R2, R3), the protein
centre-of-mass was positioned 2.183~nm from the upper interface. All systems
underwent steepest-descent energy minimisation followed by 100~ps NVT
equilibration of the solvent (protein restrained), 100~ps NVT slab equilibration
(all unrestrained), and 10~ns NPT equilibration before production.

\subsection{MD Parameters}

All simulations used the CHARMM36m force field\cite{Huang2017} with GROMACS
2020.4.\cite{Abraham2015} Key parameters are summarised in Table~\ref{tab:mdparams}.
No pressure coupling was applied, as required for slab geometry to prevent
artificial compression of the vacuum phase. Dispersion correction was omitted
in accordance with CHARMM36m slab simulation best practices.

\begin{table}[h!]
  \centering
  \caption{MD simulation parameters.}
  \label{tab:mdparams}
  \begin{tabular}{ll}
    \toprule
    Parameter & Value \\
    \midrule
    Force field           & CHARMM36m \\
    Water model           & TIP3P \\
    Temperature           & 298~K (V-rescale, $\tau = 0.1$~ps) \\
    Pressure coupling     & None (slab geometry) \\
    Electrostatics        & PME, grid spacing 0.16~nm, order 4 \\
    VdW                   & Force-switch 1.0--1.2~nm \\
    Integrator            & Leap-frog, $\Delta t = 2$~fs \\
    Bond constraints      & LINCS (all bonds) \\
    Dispersion correction & None \\
    \bottomrule
  \end{tabular}
\end{table}

\subsection{Analysis}

Trajectory analysis used MDAnalysis~2.10.0.\cite{MichaudAgrawal2011,Gowers2016}
Periodic boundary artefacts were corrected with the \texttt{unwrap}
transformation before all structural calculations. The air--water interface
position was defined as the upper boundary of the water slab, detected from
the water density profile.

\textit{Z-position and interface distance.} The protein centre-of-mass
$Z$-coordinate and the minimum distance between any protein atom and the
interface were tracked at every 100~ps frame. Adsorption was defined as a
sustained minimum distance $<$0.5~nm for $>$10 consecutive nanoseconds.

\textit{RMSD.} Root-mean-square deviation was calculated with the MDAnalysis
\texttt{RMSD} class using the initial structure as reference.
Results were extracted as \texttt{R.results.rmsd[:,2]}. Four groups were tracked:
backbone (N, C$\alpha$, C), beta-sheet strands (A--H), alpha-helix
(residues 130--142), and hydrophobic patch (calyx-lining residues).

\textit{SASA.} Solvent-accessible surface area was computed using the
\texttt{freesasa} library\cite{Mitternacht2016} interfaced through MDAnalysis.
Total, hydrophobic, hydrophilic, and calyx-patch contributions were computed
at each frame. The probe radius was set to 0.14~nm (water).

\textit{RMSF.} Per-residue root-mean-square fluctuation of C$\alpha$ atoms
was computed over the production trajectory after alignment.

\textit{Orientation.} The hydrophobic patch vector was defined as the unit
vector from the protein centre-of-mass to the calyx centroid. The angle
between this vector and the positive $Z$-axis (toward the interface) was
tracked as a function of time.

% -----------------------------------------------------------------------
\section{Results}
% -----------------------------------------------------------------------

\subsection{BLG Does Not Adsorb Within 1{,}000~ns from Bulk: Establishing the \texorpdfstring{$>$1~\textmu{}s}{>1 us} Timescale}

Figure~\ref{fig:z_center} shows the $Z$-position of the BLG centre-of-mass
and the minimum protein--interface distance over the 1{,}000~ns SET~1A
trajectory. The protein undergoes diffusive exploration of the slab
($Z$ range: 15.91--19.13~nm) without adsorbing at either interface. The
minimum distance to the upper interface was 1.598~nm, reached at $t = 570$~ns.

Adsorption was not achieved at any point during the 1{,}000~ns run under our
criterion (minimum distance $<$0.5~nm sustained for $>$10~ns). This establishes,
for the first time by atomistic simulation, that spontaneous BLG adsorption
from bulk solution requires $>$1~\textmu{}s. The result is consistent with
experimental observations of a kinetic energy barrier to BLG
adsorption\cite{Cornec1999} and with the experimental timescale of
seconds-to-minutes for measurable adsorption in tensiometry
experiments.\cite{Ulaganathan2019}

\begin{figure}[h!]
  \centering
  \includegraphics[width=0.9\columnwidth]{../results/figures/z_position/CENTER_1000ns_z_position.png}
  \caption{Z-position of the BLG centre-of-mass (top) and minimum
    protein--interface distance (bottom) over the 1{,}000~ns SET~1A
    trajectory. The dashed line indicates the adsorption threshold
    (0.5~nm). The protein approaches the interface most closely at
    $t = 570$~ns (1.598~nm) but does not adsorb within 1{,}000~ns,
    establishing a $>$1~\textmu{}s spontaneous adsorption timescale.}
  \label{fig:z_center}
\end{figure}

\subsection{Pre-Adsorption Conformational Activation: Hydrophobic SASA Spikes in Bulk}

Despite the absence of adsorption, the SET~1A trajectory reveals a striking
pattern of transient hydrophobic exposure (Figure~\ref{fig:sasa_center}).
The mean hydrophobic SASA is \SI{26.20 \pm 6.23}{\nano\metre\squared},
but three distinct activation windows are observed: $t \approx 259$~ns
(\SI{39.3}{\nano\metre\squared}), $t \approx 759$--779~ns
(\SI{39.3}--\SI{36.0}{\nano\metre\squared}), and $t \approx 940$~ns
(\SI{42.7}{\nano\metre\squared}, the largest observed). These spikes represent
a \SIrange{50}{63}{\percent} increase over the mean hydrophobic SASA.

Critically, the largest activation event ($t = 940$~ns, \SI{42.7}{\nano\metre\squared})
occurs at the end of the trajectory, when the protein is closest to the interface.
The progressive intensification of activation events --- alongside the protein's
drift toward the interface --- is consistent with a positive feedback mechanism
in which partial unfolding increases the thermodynamic driving force for
adsorption.

\begin{figure}[h!]
  \centering
  \includegraphics[width=0.9\columnwidth]{../results/figures/sasa/CENTER_1000ns_sasa.png}
  \caption{Hydrophobic solvent-accessible surface area (SASA) over the
    1{,}000~ns SET~1A trajectory. Three pre-adsorption activation windows
    are highlighted (shaded regions): $t \approx 259$~ns, 759--779~ns,
    and 940~ns. Dashed line indicates mean hydrophobic SASA
    (\SI{26.20}{\nano\metre\squared}). Spikes reaching up to
    \SI{42.7}{\nano\metre\squared} represent transient exposure of the
    hydrophobic calyx in bulk water.}
  \label{fig:sasa_center}
\end{figure}

\subsection{Structural Basis of Activation: Beta-Sheet Disruption with Alpha-Helix Stability}

RMSD analysis (Figure~\ref{fig:rmsd_center}) reveals the structural basis of
the activation events. Beta-sheet RMSD (mean \SI{0.260}{\nano\metre}) reaches
a maximum of \SI{0.502}{\nano\metre} during the activation windows,
indicating transient disruption of the barrel core. In contrast, alpha-helix
RMSD remains low throughout (mean \SI{0.101}{\nano\metre},
max \SI{0.299}{\nano\metre}) --- more than 2-fold below beta-sheet values.
The hydrophobic patch RMSD (mean \SI{0.197}{\nano\metre}) reflects calyx
mobility correlated with the SASA spikes.

This differential structural response --- beta-barrel disruption with
alpha-helix stability --- defines the activation mechanism. We propose that
the alpha-helix (residues 130--142), which flanks the calyx on one side,
acts as a structural scaffold that maintains the calyx in an accessible
orientation during insertion into the interface, while the beta-strands
transiently ``crack open'' to expose the hydrophobic interior.

\begin{figure}[h!]
  \centering
  \includegraphics[width=0.9\columnwidth]{../results/figures/rmsd/CENTER_1000ns_rmsd.png}
  \caption{RMSD of four structural groups over the SET~1A trajectory:
    backbone (black), beta-sheet strands A--H (blue), alpha-helix
    residues 130--142 (red), and hydrophobic patch (orange). The
    alpha-helix is the most stable element throughout (mean
    \SI{0.101}{\nano\metre}), while beta-sheet RMSD peaks at
    \SI{0.502}{\nano\metre} during pre-adsorption activation windows.}
  \label{fig:rmsd_center}
\end{figure}

\subsection{Near-Interface Replicas Confirm Activation and Identify Best Candidate}

The three Plan~B replicas, starting 2.183~nm from the upper interface,
independently reproduce the pre-adsorption activation behaviour
(Figure~\ref{fig:replicas}).

Replica~1 (R1) is the strongest candidate for imminent adsorption: it achieves
the closest approach of any simulation (\SI{1.35}{\nano\metre} at $t = 407.6$~ns)
and shows the most dramatic activation --- hydrophobic SASA spiking to
\SI{52.48}{\nano\metre\squared} at $t = 259$~ns, a 2.7-fold increase over the
$t = 0$ value of \SI{19.48}{\nano\metre\squared}. R1 also shows the highest
hydrophobic patch RMSD (mean \SI{0.305}{\nano\metre}), indicating maximum calyx
mobility. Replica~2 and Replica~3 approach to within 1.78 and 1.58~nm
respectively but show lower activation amplitudes, consistent with their greater
residual distance from the interface at the end of the 500~ns runs.

Across all replicas, the alpha-helix RMSD remains $\leq$\SI{0.137}{\nano\metre},
confirming that helix stability is a robust feature of BLG in bulk regardless
of proximity to the interface.

\begin{figure}[h!]
  \centering
  % Panel A: Z-position
  \includegraphics[width=0.9\columnwidth]{../results/figures/z_position/REPLICA_1_z_position.png}\\[4pt]
  \includegraphics[width=0.9\columnwidth]{../results/figures/sasa/REPLICA_1_sasa.png}\\[4pt]
  \includegraphics[width=0.9\columnwidth]{../results/figures/rmsd/REPLICA_1_rmsd.png}
  \caption{Plan~B Replica~1 results (best candidate, closest approach
    \SI{1.35}{\nano\metre} at $t = 407.6$~ns). \textit{Top:} Z-position
    and interface distance. \textit{Middle:} Hydrophobic SASA --- spike
    to \SI{52.48}{\nano\metre\squared} at $t = 259$~ns represents a
    2.7-fold pre-adsorption activation event. \textit{Bottom:} RMSD ---
    hydrophobic patch most mobile (\SI{0.305}{\nano\metre}); alpha-helix
    stable (\SI{0.111}{\nano\metre}).}
  \label{fig:replicas}
\end{figure}

\subsection{Loop~BC as Predicted First-Contact Region}

Per-residue RMSF analysis of the SET~1A trajectory reveals that Loop~BC
(residues 30--35) is the most flexible region of BLG in bulk, with
RMSF $\approx$\SI{0.53}{\nano\metre} --- 2--4-fold higher than the beta-strands
and 5-fold higher than the alpha-helix (Figure~\ref{fig:rmsf}).
This loop is located at the entrance to the hydrophobic calyx and has been
identified as structurally mobile in crystallographic studies of ligand
binding.\cite{Bello2014}

On the basis of its high flexibility and calyx-adjacent position, we predict
that Loop~BC serves as the primary contact antenna for the air--water interface.
Direct confirmation requires trajectory frames in which adsorption occurs
($<$0.5~nm) and will be reported upon completion of the extended replica
simulations.

\begin{figure}[h!]
  \centering
  % Placeholder — RMSF figure to be generated from trajectory
  % \includegraphics[width=0.9\columnwidth]{../results/figures/rmsf/CENTER_1000ns_rmsf.png}
  \fbox{\parbox{0.85\columnwidth}{\centering\vspace{1cm}
    \textbf{[FIGURE PLACEHOLDER]}\\
    RMSF per residue (C$\alpha$), SET~1A.\\
    Loop~BC (res.~30--35) highlighted.\\
    To be generated from trajectory.
    \vspace{1cm}}}
  \caption{Per-residue root-mean-square fluctuation (RMSF) of C$\alpha$
    atoms over the SET~1A trajectory. Loop~BC (residues 30--35,
    RMSF $\approx$\SI{0.53}{\nano\metre}) is the most flexible region
    and is predicted as the primary first-contact point at the air--water
    interface.}
  \label{fig:rmsf}
\end{figure}

\subsection{Orientation Analysis: Calyx Semi-Random in Bulk}

The hydrophobic patch vector shows a broad angular distribution in bulk
(mean $53.5 \pm 20.9^\circ$ relative to the $+Z$ axis,
Figure~\ref{fig:orientation}), consistent with isotropic tumbling of the
protein in solution. There is no significant preferential alignment of the
calyx toward the interface in bulk, indicating that the directional driving
force for orientation-specific adsorption acts primarily during or after
initial interfacial contact, not during long-range approach.

\begin{figure}[h!]
  \centering
  \includegraphics[width=0.9\columnwidth]{../results/figures/orientation/CENTER_1000ns_orientation.png}
  \caption{Hydrophobic patch vector angle relative to the $+Z$ axis
    (toward the air--water interface) over the SET~1A trajectory.
    The broad distribution (mean $53.5 \pm 20.9^\circ$) indicates
    semi-random orientation in bulk, with no preferential alignment
    during long-range approach.}
  \label{fig:orientation}
\end{figure}

% -----------------------------------------------------------------------
\section{Discussion}
% -----------------------------------------------------------------------

\subsection{A Pre-Adsorption Activation Mechanism for BLG}

The central finding of this work is that BLG does not passively diffuse to the
air--water interface as a rigid particle. Instead, it undergoes spontaneous,
transient partial unfolding events in bulk water that expose the hydrophobic
calyx and enable interfacial insertion. We term this process
\emph{pre-adsorption conformational activation}.

The proposed molecular sequence is: (1) bulk diffusion toward the interface on
timescales of hundreds of nanoseconds; (2) stochastic beta-barrel ``cracking''
events that transiently expose the hydrophobic calyx (SASA spikes); (3)
capture of an activated conformation at the interface --- the rate-limiting step
explaining the kinetic energy barrier; (4) post-adsorption rearrangement and
deeper unfolding (beyond the timescale of current simulations). The alpha-helix
remains intact throughout steps 1--3, acting as a structural scaffold that
maintains calyx accessibility during insertion.

This mechanism is distinct from what has been observed at oil--water interfaces.
Jordens \textit{et al.}\cite{Jordens2015} reported that BLG adsorption at
decane--water preserves secondary structure while exposing hydrophobic residues.
In that study, the protein was pre-placed at the interface and did not need to
overcome the bulk-to-interface activation barrier we observe here. Our air--water
simulations, starting from native-state BLG in bulk, capture the earlier and
energetically more demanding part of the adsorption process.

The most closely related study is Chaudhri \textit{et al.},\cite{Chaudhri2024}
who observed that small conformational changes promote monoclonal antibody
adsorption at the water/vapour interface. A key distinction is that those
authors generated partially unfolded conformations by thermal pre-stressing,
whereas our activation events arise \emph{spontaneously from the native fold}
during unbiased MD. This demonstrates that the activation is thermodynamically
accessible under physiological conditions, not a consequence of imposed
perturbation.

\subsection{Reconciling the Experimental Energy Barrier with Atomistic Simulation}

The kinetic energy barrier to BLG adsorption was established experimentally
by Cornec \textit{et al.},\cite{Cornec1999} who showed that BLG adsorption is
diffusion-controlled at short times but becomes barrier-limited at longer times.
Our simulations provide the first molecular explanation: the barrier arises
because adsorption requires a rare conformational fluctuation (the activation
event) in which the beta-barrel partially opens. The probability of this
fluctuation determines the pre-exponential factor in the Arrhenius-like
adsorption rate.

The $>$1~\textmu{}s timescale established here is a lower bound. The progressive
intensification of activation events through the 1{,}000~ns SET~1A trajectory
--- with the largest event at $t = 940$~ns --- suggests that adsorption would
occur in the 1--2~\textmu{}s range. Extended replica simulations (to 1{,}000~ns)
will refine this estimate.

\subsection{Implications for Alpha-Helix Stability and Film Rheology}

The consistent stability of the alpha-helix across all simulations has
implications for the rheology of adsorbed BLG films. Neutron reflectometry
studies\cite{Ulaganathan2019} show that adsorbed BLG films are viscoelastic,
with storage moduli consistent with a partially unfolded but not completely
denatured protein. Our data suggest that the alpha-helix survives at least
the pre-adsorption and approach phases intact. Post-adsorption SRCD studies
of BLG at oil--water interfaces report an \emph{increase} in non-native helix
content --- an observation we interpret as a consequence of interfacial
insertion (unfolding of beta-strands into helical conformations under the
influence of the interface), not a prerequisite. The helix is the stable
structural element that provides the rigidity observed in BLG interfacial
films.

\subsection{Loop~BC as First-Contact Antenna}

The identification of Loop~BC (residues 30--35) as the most flexible
surface-exposed region provides a testable prediction. If Loop~BC is indeed
the first contact point, the adsorption event should be detectable by a
sharp reduction in Loop~BC--interface distance before any other region reaches
the interface. This will be quantified once adsorption occurs in the
extended replica simulations. Experimental support comes from hydrogen--deuterium
exchange studies of BLG adsorbed at interfaces, which show that the region
around Loop~BC is among the first to become protected from exchange ---
consistent with early burial at the interface.

\subsection{Limitations}

This study does not report a complete adsorption event. The adsorption criterion
(minimum distance $<$0.5~nm sustained for $>$10~ns) was not met in any
simulation. The Plan~B replicas are being extended to 1{,}000~ns; confirmation
of adsorption and post-adsorption analysis will be reported in a follow-up
study. Additionally, only the monomeric BLG form was studied; the dimer
(dominant at neutral pH at higher concentrations) may exhibit different
activation dynamics. Finally, the use of the TIP3P water model, while
standard for CHARMM36m slab simulations, may underestimate diffusion rates
relative to more accurate water models.

% -----------------------------------------------------------------------
\section{Conclusions}
% -----------------------------------------------------------------------

We have performed the first unbiased atomistic molecular dynamics simulation
of $\beta$-Lactoglobulin at the air--water interface, revealing a
pre-adsorption conformational activation mechanism that explains decades of
experimental observations. Our key findings are:

\begin{enumerate}
  \item \textbf{Pre-adsorption conformational activation}: BLG undergoes
        spontaneous, transient partial unfolding in bulk water --- characterised
        by hydrophobic SASA spikes up to \SI{52.5}{\nano\metre\squared} --- that
        precede interfacial approach and represent the molecular prerequisite
        for adsorption.

  \item \textbf{$>$1~\textmu{}s spontaneous adsorption timescale}: The 1{,}000~ns
        CENTER simulation does not achieve adsorption, establishing the first
        atomistic lower bound on the spontaneous adsorption timescale and
        providing a molecular explanation for the experimentally observed kinetic
        energy barrier.

  \item \textbf{Differential secondary structure response}: The alpha-helix
        (residues 130--142) is structurally stable throughout all simulations
        (mean RMSD $\leq$\SI{0.137}{\nano\metre}), while the beta-barrel is
        transiently disrupted during activation (max RMSD \SI{0.502}{\nano\metre}).
        The helix acts as a scaffold maintaining calyx orientation during
        interfacial insertion.

  \item \textbf{Loop~BC as predicted first-contact region}: Per-residue RMSF
        analysis identifies Loop~BC (residues 30--35, RMSF $\approx$
        \SI{0.53}{\nano\metre}) as the most flexible surface region and
        predicted primary contact point at the air--water interface.
\end{enumerate}

These results provide a molecular foundation for understanding and engineering
BLG-stabilised foam and establish a simulation framework for Phase~2 of this
study: the contrasting adsorption mechanism of intrinsically disordered
$\beta$-Casein.

% -----------------------------------------------------------------------
\begin{acknowledgement}
[Acknowledgements to be added before submission — funding, HPC access, advisor]
\end{acknowledgement}

% -----------------------------------------------------------------------
\bibliography{references}

\end{document}
